\chapter{Frontotemporal Dementia}
\index{Frontotemporal Dementia}

\section*{Definition and Overview}
Frontotemporal dementia (FTD) is a group of progressive brain disorders that mainly affect the frontal and temporal lobes, the areas responsible for personality, behaviour, language, and decision-making. Unlike Alzheimer's disease, which usually begins with memory loss, FTD typically begins with changes in personality and behaviour or with difficulty speaking and understanding language, while memory may remain relatively intact early on. It usually begins at a younger age than other dementias, most often between 45 and 65, and it has a strong genetic component in a proportion of cases. There is no cure, but a range of strategies and medicines manage symptoms and support families. This chapter explains the types, features, diagnosis, and management of FTD.

\section*{Epidemiology}
Frontotemporal dementia is the second most common cause of early-onset dementia (before age 65), accounting for around 10--20% of dementias in this age group and up to 3% of all dementias. It affects around 1 in 100,000 people overall, and the numbers are similar in men and women. In many countries, several thousand people live with FTD. Around 30--40% of cases are familial, linked to specific gene mutations, and a family history is therefore common. The average survival after diagnosis is around 7--13 years, and the condition has a major impact on families because of the young age of onset and the behavioural changes.

\section*{Aetiology and Pathophysiology}
\begin{itemize}
    \item \textbf{Protein accumulation:} abnormal clumps of proteins build up in the brain's frontal and temporal lobes; the main proteins are tau and TDP-43
    \item \textbf{Types of FTD:} the behavioural variant (bvFTD), affecting personality and behaviour, and the language variants (primary progressive aphasias), affecting speech and comprehension
    \item \textbf{Brain changes:} the affected lobes shrink (atrophy), and the neurons that regulate social behaviour, emotion, and language are lost
    \item \textbf{Genetics:} mutations in several genes (including MAPT, GRN, and C9orf72) cause familial FTD, and the same genes are linked to motor neurone disease
    \item \textbf{Why behaviour changes:} the frontal lobes govern judgement, inhibition, empathy, and social conduct; their loss removes these controls
    \item \textbf{Why language fails:} the temporal and frontal language areas are lost in the language variants
    \item \textbf{Associated conditions:} some people with FTD also develop motor neurone disease, with muscle weakness
\end{itemize}

\section*{Clinical Features}
\begin{itemize}
    \item Behavioural variant: marked changes in personality, loss of empathy, disinhibition, inappropriate social behaviour, and poor judgement
    \item Apathy and loss of interest, or conversely impulsivity and overactivity
    \item Compulsive or repetitive behaviours, and changes in eating, such as cravings for sweet foods and overeating
    \item Language variant: difficulty finding words, naming objects, and understanding conversation, with increasingly halting speech
    \item Preserved memory and orientation early on, which can delay diagnosis
    \item Difficulty with planning, organisation, and decision-making
    \item Loss of insight: the person may not recognise the changes
    \item Physical symptoms in late stages, including rigidity and difficulty swallowing
    \item In some people, muscle weakness from associated motor neurone disease
\end{itemize}

\section*{Differential Diagnosis}
The features overlap with several conditions.
\begin{itemize}
    \item Alzheimer's disease, in which memory loss dominates
    \item Psychiatric conditions, including depression, bipolar disorder, and personality disorders, which FTD can mimic early on
    \item Vascular dementia and other dementias
    \item Motor neurone disease, which can coexist with FTD
    \item Primary psychiatric and neurological causes of apathy and disinhibition
    \item Specialist assessment, imaging, and genetic testing distinguish these
\end{itemize}

\section*{When to Seek Medical Care}
Anyone with progressive changes in personality, behaviour, or language should be assessed by a doctor, with referral to a memory clinic or neurologist. Early diagnosis matters for safety, planning, and family support.

\textbf{Contact your local emergency services for:}
\begin{itemize}
    \item Sudden confusion, weakness, or speech difficulty, which may indicate a stroke
    \item Self-harm, aggression that is dangerous, or a crisis situation
    \item Falls with head injury
    \item Difficulty breathing or swallowing emergencies
    \item Suicidal thoughts or severe distress
\end{itemize}

\section*{Diagnosis and Investigations}
\begin{itemize}
    \item \textbf{History:} detailed accounts from the person and family about changes in behaviour, personality, and language, which are essential to the diagnosis
    \item \textbf{Cognitive and behavioural assessment:} neuropsychological testing by a specialist team
    \item \textbf{Brain imaging:} MRI shows shrinkage of the frontal and temporal lobes, and PET scans can show reduced activity in these areas
    \item \textbf{Genetic testing:} for people with a family history, identifying the specific gene
    \item \textbf{Blood tests:} to exclude other causes of cognitive and behavioural change
    \item \textbf{Lumbar puncture:} in selected cases to examine spinal fluid for other dementia proteins
\end{itemize}

\section*{Management}
\begin{itemize}
    \item \textbf{Symptom management:} medicines to reduce disinhibition, agitation, and compulsive behaviour, and to treat depression and anxiety
    \item \textbf{Behavioural strategies:} structured routines, predictable environments, redirection, and avoiding confrontation
    \item \textbf{Speech therapy:} for language difficulties, including communication aids and strategies for families
    \item \textbf{Safety planning:} driving assessment, supervision where needed, and management of eating and swallowing changes
    \item \textbf{Occupational therapy:} adapting the home and daily routines
    \item \textbf{Family support:} education, counselling, respite care, and carer support, which are essential in FTD
    \item \textbf{Legal and financial planning:} early planning for guardianship, powers of attorney, and finances while the person can participate
    \item \textbf{Specialist services:} memory clinics, dementia teams, and community supports
    \item \textbf{Later care:} increasing assistance with daily living, and palliative care at the end of life
\end{itemize}

\section*{Complications}
\begin{itemize}
    \item Behavioural problems that strain relationships and require professional management
    \item Financial and legal difficulties from impaired judgement
    \item Falls and injuries from motor and balance changes
    \item Swallowing difficulties with choking and aspiration pneumonia
    \item Carer burnout and family breakdown
    \item Shortened life expectancy, especially with associated motor neurone disease
\end{itemize}

\section*{Prognosis and Outcomes}
Frontotemporal dementia is a progressive condition without a cure, and survival averages 7--13 years from diagnosis. However, the course varies, and good care significantly improves quality of life for the person and the family. Behavioural symptoms respond to structured management and medicines in many cases, and early planning prevents financial and safety crises. Families who access education, respite, and counselling cope better, and support services are central to care. Research into the protein abnormalities and genetic forms is advancing treatment prospects.

\section*{Prevention and Self-Care}
\begin{itemize}
    \item Seek specialist assessment early for personality, behaviour, or language changes
    \item Request genetic counselling if there is a family history of early-onset dementia or motor neurone disease
    \item Establish a routine and a calm, predictable environment
    \item Plan early for driving, finances, legal matters, and future care
    \item Access dementia support services, respite, and carer counselling
    \item Join a dementia-specific support group, ideally for FTD
    \item Look after your health and wellbeing as a carer
    \item Stay informed about research and clinical trials
\end{itemize}

\section*{Sources and Further Reading}
The following reputable sources provide further information for healthcare students and patients seeking reliable, current advice about frontotemporal dementia.
\begin{itemize}
    \item Dementia Australia. \textit{Frontotemporal dementia.} https://www.dementia.org.au/
    \item Frontotemporal Dementia Support Group Australia. \textit{Support and information.} https://www.ftdsupport.com.au/
    \item Healthdirect Australia. \textit{Frontotemporal dementia.} https://www.healthdirect.gov.au/
    \item NHS. \textit{Frontotemporal dementia.} https://www.nhs.uk/conditions/frontotemporal-dementia/
    \item Association for Frontotemporal Degeneration. \textit{Resources.} https://www.theaftd.org/
    \item World Health Organization. \textit{Dementia.} https://www.who.int/
\end{itemize}
